\section{Resplendence, definability, and atomic compactness}

\begin{annexitem}{annex:p03-resplendence-padding}{thm:p03-resplendence-padding}{A decidable syntactic range with unchanged realizations}
\begin{proof}
Run an enumerator for the computably enumerable set \(p\) in discrete
steps, with at most one output per step. At a step \(s\) with output
\(\theta_s(\bar x,\bar a)\), output the formula
\[
 \theta_s(\bar x,\bar a)\land\tau_s,
 \qquad \tau_s:=\underbrace{(z=z)\land\cdots\land(z=z)}_{s+1\text{ copies}},
\]
where the displayed conjunction has a fixed right-associated syntax and
\(z\) is the first variable of \(\bar x\); if the block is empty use
\(\forall z(z=z)\) instead of \(z=z\). Do not simplify the formula.
To decide membership in the padded set, parse the formula as exactly the
displayed outer conjunction, check the canonical shape of \(\tau_s\),
read its unique length \(s+1\), and simulate the original enumerator
through step \(s\). Membership holds precisely if its step-\(s\)
output was the parsed left conjunct. Invalid syntax is rejected. Thus the
set is recursive even when the original enumerator repeats or stops.
Every \(\tau_s\) is valid in a nonempty structure, so the padded set
has exactly the original realizations and finite satisfiability. The
parameter footprint and the finite free-variable block are unchanged.
\end{proof}
\end{annexitem}

\begin{annexitem}{annex:p03-resplendence-diagram}{thm:p03-resplendence-diagram}{Finite diagram fragments can be named inside a theory model}
\begin{proof}
One direction is inclusion. For the other choose a model of
\(S\cup\operatorname{Th}(M,\bar a)\). A finite part of the full
elementary diagram mentions only finitely many further names \(\bar c\).
Its conjunction has the form \(\gamma(\bar c,\bar a)\), with a base
\(L\)-formula \(\gamma\), after all repeated names are represented by
the same variable. Since \(M\models\gamma(\bar c,\bar a)\), the
sentence \(\exists\bar z\gamma(\bar z,\bar a)\) belongs to
\(\operatorname{Th}(M,\bar a)\). Interpret these finitely many further
names by such witnesses in the chosen expanded model. It satisfies the
finite diagram fragment and every member of \(S\). Compactness gives
a model of \(S\cup\operatorname{ElDiag}(M)\). The map sending each
\(a\) to its named interpretation is elementary by
\cref{thm:diagram-tests}; inequality instances make it injective. Nothing
in this argument prevents additional unnamed elements.
\end{proof}
\end{annexitem}

\begin{annexitem}{annex:p03-resplendence-witness}{thm:p03-resplendence-witness}{Equality patterns, base consequences, and the fixed-model tuple}
\begin{proof}
Collect every named member occurring in \(S\) or \(p\) into a finite
tuple \(\bar a\) without duplicates. Consider the finitely many equality
patterns for \(\bar x\) relative to these members: each coordinate equals
one member of \(\bar a\), or belongs to a fresh equivalence class of
coordinates distinct from \(\bar a\) and from the other fresh classes.
These patterns exhaust all tuples. Since the displayed theory is
consistent, at least one pattern \(\pi\) remains consistent with it.
Substitute the prescribed parameters and one fresh variable \(y_j\)
for each fresh class into \(p\). Adjoin all inequalities
\(y_j\ne a_i\) and \(y_j\ne y_k\) for distinct classes, obtaining
a computably enumerable type \(p_\pi(\bar y,\bar a)\). If there are
no fresh classes, the resulting actual parameter tuple already works.

Otherwise define \(q(\bar y,\bar a)\) as all base \(L(\bar a)\)
formulas which follow from \(S\cup p_\pi(\bar y,\bar a)\). More
formally replace \(\bar y\) by a fixed fresh tuple of constants, enumerate
finite proofs from \(S\) and the enumerator of \(p_\pi\), retain base
formulas in these constants, and replace the constants back by variables.
Finite-proof enumeration is effective even for computably enumerable
axioms: dovetail proof candidates with the enumerators and certify each
nonlogical axiom when it appears. No diagram axiom is used in this
enumeration. Hence \(q\) is computably enumerable over finite \(\bar a\).

For every finite \(q_0\subseteq q\), a model of
\(\operatorname{ElDiag}(M)\cup S\cup p_\pi\) contains a tuple
satisfying \(q_0\). Its elementary copy of \(M\) therefore satisfies
\(\exists\bar y\bigwedge q_0\), which is a base sentence with those
parameters. Thus \(q\) is finitely satisfiable in \(M\). Recursive
saturation and the padding lemma supply \(\bar b\in M\) realizing
\(q\). Since the inequalities belong to \(q\), its coordinates are
distinct and outside \(\bar a\).

Suppose \(\operatorname{ElDiag}(M)\cup S\cup p_\pi(\bar b,\bar a)\)
were inconsistent. A finite inconsistent subset uses a finite diagram
conjunction \(\theta(\bar b,\bar c,\bar a)\), true in \(M\), and
finitely many members of \(S\) and \(p_\pi\). Replace the additional
names \(\bar c\) by variables \(\bar z\). Because none of these
names occur in \(S\) or \(p_\pi(\bar b,\bar a)\), deduction and fresh
constant generalization give
\[
 S\cup p_\pi(\bar b,\bar a)\ \vdash\
             \neg\exists\bar z\theta(\bar b,\bar z,\bar a).
\]
The names for the distinct \(\bar b\) are likewise absent from \(S\)
and the uninstantiated type. Replace them uniformly by \(\bar y\) in
that finite proof. It shows
\(\neg\exists\bar z\theta(\bar y,\bar z,\bar a)\in q\).
But \(\bar b\) realizes \(q\), contradicting the actual witnesses
\(\bar c\) for \(\theta\) in \(M\). This proves consistency after
the fixed-model substitution. Recover the original tuple \(\bar x\)
from its pattern, parameters and \(\bar b\), obtaining the required
actual tuple in \(M\). The argument separates fresh names from fixed
parameters; it does not illegally generalize a parameter already in \(S\).
\end{proof}
\end{annexitem}

\begin{annexitem}{annex:p03-resplendence-countable}{thm:p03-resplendence-countable}{A fixed-domain Henkin construction with a type schedule}
\begin{proof}
Let \(D=\operatorname{ElDiag}(M)\). The diagram lemma gives consistency
of \(D\cup T\). Because \(M\) and \(L^+\) are countable, fix external
enumerations of: all \(L^+(M)\)-sentences; all one-free-variable
\(L^+(M)\)-formulas; and all recursive types in finite variable blocks
over finite tuples from \(M\). Interleave these three countable lists.
The last is a set-theoretic enumeration of the total deciders and parameter
tuples, not an effective decision of which machine indices are total.

Construct \(S_0=T,S_1,S_2,\ldots\). Each \(S_i\) is computably
enumerable, mentions only finitely many members of \(M\), extends its
predecessor, and \(D\cup S_i\) is consistent. For a sentence task
\(\sigma\), adjoin \(\sigma\) if consistent with \(D\cup S_i\),
otherwise adjoin \(\neg\sigma\). At least one choice is consistent,
and in the second case every model of \(D\cup S_i\) satisfies the
negation. For a witness task \(\varphi(x)\), if
\(D\cup S_i\cup\{\exists x\varphi(x)\}\) is consistent, the witness
lemma supplies \(b\in M\) for which adjoining \(\varphi(b)\) is
consistent. Otherwise adjoin \(\neg\exists x\varphi(x)\).

For a recursive-type task \(p(\bar x,\bar a)\), if
\(D\cup S_i\cup p(\bar d,\bar a)\) is consistent, use the witness
lemma and adjoin the entire \(p(\bar b,\bar a)\) for an actual
\(\bar b\in M\). This is a computably enumerable axiom set with
finite parameter footprint. If it is inconsistent, compactness gives a
finite \(p_0\subseteq p\) such that
\[
 D\cup S_i\models\neg\exists\bar x\bigwedge p_0(\bar x,\bar a).
\]
Indeed a finite inconsistent fragment may include finite parts of \(D\)
and \(S_i\); these belong to the entire \(D\cup S_i\), and the new
\(\bar d\) names are absent there. Adjoin this negated existential
sentence. At each stage only one finite-footprint type and finitely many
sentences are added, so all hypotheses of the witness lemma remain valid.
The choices are semantic choices; no computable expansion is asserted.

Put \(S_\infty=\bigcup_iS_i\). Compactness gives a model \(K\) of
\(D\cup S_\infty\), and let \(h:M\to K\) be its elementary base
embedding. Every existential \(L^+\)-formula with parameters in
\(h[M]\) true in \(K\) has a witness in \(h[M]\): its scheduled task
could not have taken the negative branch, and its positive branch adjoined
the named witness. This includes \(\exists x(x=c)\) for each new
constant and \(\exists x(x=f(\bar a))\) for each new function value.
Thus \(h[M]\) contains all new constants and is closed under all new
functions. With the induced relations it is an \(L^+\)-substructure.

It is elementary in \(K\), not merely a substructure. Induct on formulas
with parameters from \(h[M]\). Atomic formulas agree by the induced
substructure; Boolean clauses preserve agreement. In the existential
clause a witness in the substructure gives one in \(K\), and a witness
in \(K\) has a scheduled named witness in the substructure. Apply the
induction hypothesis to its matrix. Universal clauses follow by negation.
Transport this \(L^+\)-structure along \(h\) to the original domain
of \(M\). In particular new functions are
\(h^{-1}(f^K(h\bar a))\), and new relations hold exactly when they
hold of \(h\bar a\) in \(K\). The old interpretations agree with
those of \(M\), the named parameters remain their actual values, and
the resulting \(M^+\) satisfies \(T\).

Let a recursive expanded type \(p(\bar x,\bar a)\) be finitely
satisfiable in \(M^+\). At its scheduled stage every finite part of
\(D\cup S_i\cup p(\bar d,\bar a)\) has a model: use \(M^+\)
and a witness for the finitely many type formulas. It was therefore
consistent, so that task adjoined \(p(\bar b,\bar a)\). The expansion
just constructed satisfies this whole type at \(\bar b\). Hence it
is recursively saturated. This supplies the strengthened assertion by
an actual additional schedule, beyond the single-witness proof printed
in the source.

Conversely suppose the recursive-theory expansion property holds. For a
recursive finitely satisfiable base type \(p(\bar x,\bar a)\), add
fresh constants \(\bar d\) and take \(T=p(\bar d,\bar a)\).
Substitution of these fresh names gives a recursive set of sentences.
Its union with \(\operatorname{Th}(M,\bar a)\) is consistent by
compactness, using actual witnesses for each finite part. An expansion
of the fixed \(M\) satisfying \(T\) interprets \(\bar d\) by a
realization in \(M\). This proves recursive saturation. Finally a
single sentence with finitely many new symbols is a recursive theory in
a recursive extension, so the already proved forward assertion supplies
finite-sentence resplendence too.
\end{proof}
\end{annexitem}

\begin{annexitem}{annex:p03-resplendence-extension}{thm:p03-resplendence-extension}{Two countable rounds of actual elementary extensions}
\begin{proof}
Start with \(M_0=M\). At outer stage \(n\), enumerate all recursive
types in finite variable blocks over finite tuples from \(M_n\) which
are finitely satisfiable in \(M_n\). There are countably many: a type
has a standard machine index and a finite parameter tuple. This enumeration
need not be computable; finite satisfiability is not being decided by an
algorithm. Build an inner elementary chain \(K_0=M_n,K_1,\ldots\),
realizing one of these types \(p_j\) at step \(j+1\).

The theory \(\operatorname{ElDiag}(K_j)\cup p_j(\bar d)\) is finitely
satisfiable. A finite part of \(p_j\) has an actual witness in \(M_n\)
and therefore in \(K_j\); with that choice for \(\bar d\), the
structure \(K_j\) satisfies every selected diagram sentence as well.
Compactness gives an elementary extension with a realization. Apply
downward Lowenheim--Skolem, including the elementary image of \(K_j\)
and that finite tuple in the chosen subset, to obtain a countable elementary
substructure. Rename its embedded copy of \(K_j\) to make this an actual
extension, and call it \(K_{j+1}\). The diagram test and the elementary
substructure step preserve every old formula. Put
\(M_{n+1}=\bigcup_jK_j\); the elementary-chain theorem makes it a
countable elementary extension of \(M_n\) realizing the whole list.

Put \(N=\bigcup_nM_n\). It is countable and \(M\preccurlyeq N\).
If a recursive \(p(\bar x,\bar a)\) is finitely satisfiable in \(N\),
its finite parameter tuple lies in some \(M_n\). For each finite
\(p_0\subseteq p\), the sentence
\(\exists\bar x\bigwedge p_0(\bar x,\bar a)\) true in \(N\)
is true in \(M_n\) by elementarity. Thus the entire type was on that
stage's list and was realized in \(M_{n+1}\). Its realization remains
one in \(N\). Hence \(N\) is recursively saturated. This proof keeps
the parameters in one outer stage before scheduling their type; a single
enumeration over \(M_0\) would not suffice. The resplendence theorem
now applies to \(N\), not retrospectively to \(M_0\).
\end{proof}
\end{annexitem}

\begin{annexitem}{annex:p03-resplendence-homogeneity}{thm:p03-resplendence-homogeneity}{A recursive comparison type and the full back-and-forth}
\begin{proof}
For tuples of equal complete type and \(c\in M\), enumerate all base
formulas \(\varphi(\bar y,x)\) and form the computably enumerable type
\[
 q(x;\bar a,\bar b,c)=
 \{\varphi(\bar a,c)\longrightarrow\varphi(\bar b,x):
                                      \varphi(\bar y,x)\text{ an }L\text{-formula}\}.
\]
For a finite part keep only the matrices whose premises are true at
\((\bar a,c)\). The sentence asserting a witness for their conjunction
is true at \(\bar a\), hence true at \(\bar b\), since those tuples
have the same complete type. Its witness satisfies all the finite
implications, including those with false premises. The padding lemma
and recursive saturation give a realization \(d\). Both \(\varphi\)
and \(\neg\varphi\) occur in the enumeration, so truth of every
\(\varphi(\bar a,c)\) agrees with truth of
\(\varphi(\bar b,d)\). No recursive membership test for
\(\operatorname{tp}(\bar a,c)\) was assumed. If \(c\) is outside the
source tuple, the corresponding inequality instances force \(d\)
outside its image; thus the enlarged finite map is injective and elementary.

For a countable \(M\), enumerate its elements. Starting with the given
finite elementary map, at an even step take the first element not in its
domain and extend by the just proved construction. At an odd step apply
the construction to the inverse finite elementary map, adding the first
element outside its range, then invert back. Skip a step if the chosen
element is already included. The union has every element in its domain
and range, is bijective, and preserves every formula on every finite
tuple because that tuple occurs at a finite stage. It therefore commutes
with all functions and preserves and reflects all relations, and is an
automorphism extending the initial map. For finite \(M\), the same
process terminates with a total bijection.
\end{proof}
\end{annexitem}

\begin{annexitem}{annex:p03-resplendence-definability}{thm:p03-resplendence-definability}{Mixed type cells or a finite defining formula}
\begin{proof}
A relation defined by a base formula with parameters \(\bar a\) is
preserved by every base automorphism fixing them, by induction on that
formula. For the converse, the nullary case is already defined by truth
or falsity, so suppose \(n>0\). Assume \(R\) is not definable over
\(\bar a\). For any finite list of base formulas
\(\varphi_1(\bar x,\bar a),\ldots,\varphi_k(\bar x,\bar a)\),
partition \(M^n\) into the finitely many signed truth-pattern cells
\(\theta_\sigma(\bar x,\bar a)\), \(\sigma\in\{0,1\}^k\).
If each cell had constant \(R\)-membership, the disjunction of exactly
the cells contained in \(R\) would be a base defining formula. Therefore
some cell has both an \(R\)-tuple and a non-\(R\)-tuple.

Consequently the computably enumerable type in the \emph{expanded} language
\[
 \{R(\bar x),\neg R(\bar y)\}\ \cup\
 \{\varphi(\bar x,\bar a)\leftrightarrow\varphi(\bar y,\bar a):
                                       \varphi\text{ a base formula}\}
\]
is finitely satisfiable in \((M,R)\). Padding and the explicitly assumed
recursive saturation of this given expansion yield actual tuples
\(\bar b,\bar c\) with equal complete base type over \(\bar a\),
but opposite membership. The base reduct is recursively saturated because
every recursive base type finitely satisfiable there is also finitely
satisfiable in the expansion. The previous back-and-forth theorem extends
the finite elementary map \((\bar a,\bar b)\mapsto(\bar a,\bar c)\)
to a base automorphism fixing \(\bar a\). It fails to preserve \(R\),
contradicting the asserted invariance. Thus \(R\) is definable.

This also makes the finite formula construction explicit: if no separating
automorphism exists, the displayed type cannot be finitely satisfiable,
so finitely many \(\varphi_i\) suffice to make each truth-pattern cell
uniform in membership. Select the positive cells and form their finite
disjunction, using falsity for no such cell. The selecting decision is
external, not an effective definability algorithm. Its resulting formula
contains only \(L\) and the named finite parameters, not \(R\).
\end{proof}
\end{annexitem}

\begin{annexitem}{annex:p03-resplendence-positive-maps}{thm:p03-resplendence-positive-maps}{The existential and universal homomorphism clauses}
\begin{proof}
First fix \(C,s,t_1,t_2\) as in the displayed satisfaction bridge.
Interpret each term recursively: a variable takes its assigned value,
a parameter name takes its named member of \(C\), and a function term
applies the interpreted function to its argument values. Delhomme's
native identity atom \(t_1\equiv t_2\) is true precisely when these
two resulting objects are identical. The host logical equality clause
for \(t_1=t_2\) is true precisely under the same condition. If the
objects are identical, both clauses hold; if they are different, neither
holds. This proves both directions of the bridge for every assignment,
including named parameters. It compares satisfaction, not source syntax:
the native atom in the atomic presentation remains \(\equiv\).

Term induction shows \(r(t^B(\bar b))=t^A(r\bar b)\): variables are
the assignment clause, constants are preservation of constants, and a
function application uses commutation with that function and the induction
hypotheses for its arguments. A true native identity atom has identical
term values in \(B\); applying \(r\) keeps those values identical,
so their image terms satisfy the native identity atom in \(A\).
The bridge also gives the corresponding host equality clause. Basic
relations are preserved by the homomorphism assumption. These give the
atomic induction. A conjunction retains both true conjuncts, and a
disjunction retains its true alternative. For an existential positive
formula choose its witness in \(B\); the induction hypothesis applied
to that witness makes its image a witness in \(A\).

For the universal clause assume \(r\) is onto. If
\(B\models\forall y\psi(y,\bar b)\), take any \(a\in A\) and a
preimage \(b\in B\). The instance \(\psi(b,\bar b)\) is true in
\(B\), so its image is true in \(A\) by induction. Since \(a\)
was arbitrary, the universal formula holds in \(A\). This proves the
two preservation assertions at their distinct scopes. For a retraction,
all named parameters from \(A\) are fixed and surjectivity follows from
\(r(a)=a\). Applying the induction formula by formula to a system's
assignment in \(B\) gives its returned assignment in \(A\).
\end{proof}
\end{annexitem}

\begin{annexitem}{annex:p03-resplendence-atomic-retractions}{thm:p03-resplendence-atomic-retractions}{The entire homomorphism diagram and the return of witnesses}
\begin{proof}
Suppose \(A\) is atomic compact and \(A\subseteq B\) is pure.
Introduce an unknown \(x_b\) for every \(b\in B\). Write the following
atomic system with parameters in \(A\):
\[
\begin{gathered}
 x_a\equiv a\quad(a\in A),\qquad x_{c^B}\equiv c\quad(c\text{ a constant}),\\
 f(x_{b_1},\ldots,x_{b_k})\equiv x_{f^B(b_1,\ldots,b_k)}
                          \quad(f\text{ a function},\ \bar b\in B^k),\\
 R(x_{b_1},\ldots,x_{b_k})\quad
                          (B\models R(\bar b)).
\end{gathered}
\]
In the identity atoms, \(a\) denotes its \(L(A)\)-parameter name
and the \(x_b\) are object variables, so both operands are same-sort
terms; \(\equiv\) is the native logical identity, not elementary
equivalence of structures. The relation clause includes nullary true
relations. Every finite subsystem is realized in \(B\) by the
assignment \(x_b\mapsto b\). Replace unknowns whose
indices lie in \(A\) by their parameter values, and existentially bind
the remaining finitely many unknowns. This is a pp formula with parameters
in \(A\), so purity supplies an assignment satisfying that finite
subsystem in \(A\). Atomic compactness solves the whole system. Put
\(r(b)=x_b\). The displayed clauses respectively say that \(r\) fixes
\(A\), preserves constants, commutes with every function and preserves
every relation. Thus \(r\) is the required actual retraction.

An elementary extension is pure: existential conjunctions are first-order
formulas with parameters in \(A\), and elementarity reflects their truth.
Hence the retraction property for pure extensions implies it for elementary
extensions. Suppose now that every elementary extension retracts to
\(A\), and let \(\Phi(X)\) be a finitely solvable positive system.
For each unknown adjoin a fresh constant. The theory
\(\operatorname{ElDiag}(A)\cup\Phi\) is finitely satisfiable: use
\(A\) with a solution of the finitely many system members, and assign
all other new constants an arbitrary fixed member of \(A\).
Compactness gives a model containing an elementary copy of \(A\) and
a simultaneous solution. Identify that copy with \(A\). The stipulated
retraction and the preceding positive-preservation lemma return the whole
solution, fixing every parameter. Thus \(A\) is positive compact.
Positive compactness includes atomic compactness, completing the cycle
of implications. This proof supplies the map rather than calling the
diagram's consistency a fixed-domain solution. It makes no injectivity
or elementarity claim about the retraction itself.
\end{proof}
\end{annexitem}

\begin{annexitem}{annex:p03-resplendence-topological}{thm:p03-resplendence-topological}{Closed constraints in the full assignment space}
\begin{proof}
For completeness record the product compactness step at the needed
arbitrary-set scope. A proper filter on subsets of a set extends to an
ultrafilter in the assumed metatheory: order the containing proper filters
by inclusion, observe that the union of a chain is proper, and apply Zorn's
lemma. A maximal proper filter decides every subset \(E\), since if
\(E\) cannot be added, some member of the filter misses \(E\), forcing
its complement into the filter; if it can be added, maximality already
puts it in the filter. This is the ultrafilter property.

Let compact Hausdorff spaces \(K_j\), \(j\in J\), be nonempty. Under
AC their product is nonempty. For an ultrafilter \(U\) on the product,
its coordinate image \(U_j\) is an ultrafilter on \(K_j\). The closures
of its members have the finite intersection property: a finite intersection
of members is nonempty and lies in the intersection of their closures.
Compactness supplies a point \(a_j\) in all these closures. Every open
neighborhood of \(a_j\) is in \(U_j\), since otherwise its closed
complement would be a member whose closure misses \(a_j\).
Hausdorff separation makes this point unique. The tuple \(\bar a\)
is consequently defined coordinatewise. A basic neighborhood of it restricts
only finitely many coordinates, and its cylinder sets are members of
\(U\); their intersection is a member. Thus \(U\) converges to
\(\bar a\) in the product topology. Any closed member of \(U\)
contains that limit, or its open complement would be a second, disjoint
member. A family of closed sets with the finite intersection property
generates a proper filter; extend it as above and use its limit to obtain
a common point. This proves the required product compactness, including
the empty product.

For an atomic system \(\Phi(X)\), the solution set of each member is
closed in \(A^X\). Each term is a continuous function of its finitely
many coordinates, by term induction from continuous basic functions;
a parameter is a constant function. A relation constraint is the inverse
image of its closed basic relation. A native identity constraint
\(t_1\equiv t_2\) is the inverse image of the closed diagonal by the
proved satisfaction bridge: its two term values must be the same object.
That diagonal is closed because \(A\) is Hausdorff.
Nullary constraints are either the whole assignment space or the empty
set. Finite solvability says that these closed constraint sets have the
finite intersection property. The product compactness just proved gives
an assignment in their entire intersection, proving atomic compactness.

A finite discrete space is compact Hausdorff and every operation and
relation on it is continuous or closed as required. For
\(\{0,1\}^{\mathbb N}\), the preceding product argument first makes
the carrier compact Hausdorff. Boolean complementation and the binary
Boolean operations are continuous: the inverse image of each basic cylinder
is determined by finitely many input coordinates and is open. Its equality
diagonal is closed and it has no additional basic relations. Therefore the
coordinatewise Boolean algebra has the required compatible topology and
is atomic compact. No converse from atomic compactness to the existence
of such a topology is asserted.
\end{proof}
\end{annexitem}

\begin{annexitem}{annex:p03-resplendence-atomic-powers}{thm:p03-resplendence-atomic-powers}{The filter quotient and the finite-subsystem cone solution}
\begin{proof}
For a proper filter \(F\), equality on an \(F\)-large set is an
equivalence relation: reflexivity and symmetry are immediate, and
transitivity follows by intersecting two equality sets. If finitely many
representatives are changed, they agree with their replacements on an
\(F\)-large common intersection. Function values therefore agree there,
and a basic relation's old and new truth sets agree there. Two subsets
agreeing on a filter member are in the filter simultaneously. This proves
all interpretations in the quotient well defined. Term induction gives
coordinate evaluation modulo this equivalence.

For every pp formula \(\psi\) and tuple of representing functions
\(\bar u\), we have
\[
 A^I/F\models\psi([\bar u])
 \quad\Longleftrightarrow\quad
 \{i:A\models\psi(\bar u(i))\}\in F.
\]
For a native identity atom, quotient satisfaction means that its two term
values are the identical equivalence class. By the definition of
\(\sim_F\), this holds exactly when their coordinate identity set is
in \(F\). For a basic relation atom, the required condition is the
relation definition. These give the atomic case; conjunction
uses closure under finite intersections and upward closure, which make
\(E\cap H\in F\) equivalent to both \(E,H\in F\). For the
existential case, a representing witness function in the quotient gives,
by induction, a filter member on which its matrix holds. This member is
contained in the coordinate existential truth set. Conversely, if that
truth set is in \(F\), choose an actual matrix witness at each of its
indices, and use a fixed arbitrary element of \(A\) elsewhere. AC supplies
the witness function; its matrix truth set contains the given member and
is in \(F\). Apply induction to obtain a quotient witness. Finite
conjunctions and successive existentials complete the pp induction.

For diagonal parameter tuples, the coordinate pp truth set is either
\(I\) or \(\varnothing\). Since the filter is proper, this proves
preservation and reflection of pp formulas. The diagonal is an embedding
and a pure extension. Thus atomic compactness and the preceding retraction
theorem give retractions of every reduced power, in particular every
ultrapower.

Conversely let \(\Phi(X)\) be finitely solvable and set
\(I=[\Phi]^{<\omega}\). For every finite \(u\subseteq\Phi\),
choose an assignment \(s_u:X\to A\) satisfying \(u\), extending a
finite assignment by a fixed default value if necessary. AC chooses this
family. For each unknown put \(f_x(u)=s_u(x)\). The sets
\(C_v=\{u:v\subseteq u\}\), for finite \(v\subseteq\Phi\),
generate a proper filter, since finite intersections are cones for the
finite union and are nonempty. For every \(\varphi\in\Phi\), its
coordinate truth set under \(\bar f\), with its fixed parameters given
by diagonal functions, contains \(C_{\{\varphi\}}\). The quotient
therefore realizes the entire atomic system. If every reduced power
retracts to \(A\), apply that retraction to the representing classes
to obtain a solution in \(A\). If instead only every ultrapower is
assumed to retract, extend the cone filter to an ultrafilter by the
ultrafilter construction in the preceding proof; the same truth sets
remain large and the same return argument applies. This proves both
converses, including an empty system (whose finite-subsystem index set
still contains \(\varnothing\)). No disjunction or negation clause
for arbitrary reduced powers was used.
\end{proof}
\end{annexitem}

\begin{annexitem}{annex:p03-resplendence-examples}{thm:p03-resplendence-examples}{The expansions and both separating models}
\begin{proof}
First, in an infinite equality-only structure the truth of a formula on
a finite tuple with finite named parameters depends only on its equality
pattern. This follows by induction: atomic equality tests that pattern;
Boolean combinations retain dependence; an existential variable can equal
one of the finitely many existing distinct values or be a fresh value.
All these extension patterns are available in every infinite carrier,
so the existential truth depends only on the old pattern. Negation gives
the universal clause. For fixed finite variables and parameters there
are finitely many patterns. If a type is finitely satisfiable, the sets
of patterns satisfying its members have the finite intersection property
in this finite set. Hence one pattern satisfies every member, and it has
an actual realization in \(\mathbb N\). This proves recursive saturation,
indeed realization of every type over finite parameters.

The displayed \(f\) exchanges the two members of each pair
\(\{2k,2k+1\}\); applying it twice is the identity and it has no
fixed point. Exactly one member of each pair is in \(P\), proving all
three expansion axioms directly. For any finite parameter set, choose
an even and an odd number outside it; their transposition is an automorphism
of the equality reduct fixing the parameters, but not \(P\). Thus
\(P\) is not definable there. Choose a pair \(x,f(x)\) outside that
finite set and a third unused value. A permutation fixing \(x\) and
the parameters but moving \(f(x)\) fails to preserve the function's
graph. The graph therefore is not definable either. These arguments do
not prevent new, externally chosen symbols from making the expansion.

Second, in the named equality model, every finite subset of
\(\{x\ne c_k:k\in\mathbb N\}\) is realized by an unnamed-for-that-
finite-subset number, but no number realizes the whole recursive type.
The complete theory plus the type of a fresh constant is finitely
satisfiable, using the same domain for each finite subset, so compactness
gives a model of it. The interpretation cannot be on the fixed named
domain, where every element equals some \(c_k\). This proves failure
of both recursive saturation and the recursive-theory expansion property.

An atomic system in this language consists only of equations between
unknowns and named members, or between two unknowns. Generate the
equivalence relation on the unknowns and named members by these equations.
Finite solvability prevents two different named numbers from lying in
one class: an alleged connection is a finite path of equations, already
a finite inconsistent subsystem. Assign each class containing a named
number that number, and every other class the default number \(0\).
This solves the entire system, so the structure is atomic compact.
All its automorphisms fix every number because they fix all \(c_k\).
Yet a formula with finite parameters mentions only finitely many of
these constants. Every permutation of the remaining infinite domain is
an automorphism of the finite-signature reduct actually used by that
formula. Therefore its unary solution set is constant on that complement,
and is finite or cofinite. An infinite coinfinite subset is not definable,
although invariant under every automorphism of the full named structure.
The reduct and any relation expansion still omit the displayed recursive
type; the missing saturation hypothesis is concrete.

Third, apply the elementary-extension construction to the recursive
finite-language structure \((\mathbb N,0,S,<)\), and then the expansion
theorem to its countable recursively saturated extension \(M\).
Elementarity retains strict linear order and
\(\forall y\exists z(y<z)\). A finite set of parameters has a maximum
and then a larger member, so every finite part of \(\{a<x:a\in M\}\)
has a solution. A simultaneous solution \(b\) would have to satisfy
\(b<b\), using its own parameter instance, contradicting irreflexivity.
This system uses all parameters, not a recursive type over a fixed finite
tuple; hence it contradicts atomic compactness but not recursive saturation
or the proved countable resplendence theorem.
\end{proof}
\end{annexitem}
